% Getting started with LaTeX - the whole document, no hidden parts.
% Compile me anywhere: Overleaf, TeX Live, MiKTeX.
\documentclass[11pt]{article}

\usepackage[a4paper, margin=2.5cm]{geometry}
\usepackage{amsmath}
\usepackage{booktabs}

\title{Getting Started with \LaTeX}
\author{Ken Reid \\ \texttt{kenreid.co.uk}}
\date{July 2026}

\begin{document}
\maketitle

\section{How this works}
You are reading a PDF that was generated from about forty lines of
plain text. Headings are numbered automatically, spacing and fonts
follow a professional standard, and nothing on this page was nudged
into place by hand. Change the text, recompile, and the layout takes
care of itself.

\section{Mathematics, the party trick}
Notation that word processors fight you over is native here:
\[
  \hat{\theta} \;=\; \arg\max_{\theta} \sum_{i=1}^{n} \log p(x_i \mid \theta)
\]
Inline mathematics like $e^{i\pi} + 1 = 0$ sits comfortably in a
sentence, and equation numbering, when you want it, is automatic.

\section{Structure for free}
Lists and tables stay consistent across a whole thesis:
\begin{itemize}
  \item Sections, figures, and equations number themselves.
  \item Cross-references update when things move.
  \item Bibliographies format themselves from a reference file.
\end{itemize}

\begin{center}
\begin{tabular}{lcc}
  \toprule
  Tool     & Learning curve & Ceiling \\
  \midrule
  Word     & flat           & low     \\
  Markdown & flat           & medium  \\
  \LaTeX{} & steep          & none    \\
  \bottomrule
\end{tabular}
\end{center}

\section{Your move}
Open this file in Overleaf, change something, and press recompile.
That is the entire barrier to entry.

\end{document}
